% ============================================================
% SECTION 2 — Imam al-Hassan's Pacifism (Chapter 2, Sessions 1-2)
% ============================================================
\section{Chapter 2 — al-Hassan's Pacifism}

\chapterdivider{Chapter 2}{Imam al-Hassan's Pacifism}

\begin{frame}{Not a Question of Temperament}
\bigquote{Different temperaments cannot explain why two brothers, raised the same way, chose opposite methods.}{Mutahhari's opening move, paraphrased}
\begin{itemize}[itemsep=5pt]
  \item The lazy answer: al-Hassan was peaceful, al-Husayn was brave
  \item Mutahhari's answer: both were reasoned applications of the \emph{same} jurisprudence to \emph{different} circumstances
  \item Proof by precedent: even the Prophet alternated war and peace (Hudaybiyyah vs. later battles) depending on circumstance, not temperament
\end{itemize}
\end{frame}

\begin{frame}{Categories of Jihad}
\centering\scriptsize
\renewcommand{\arraystretch}{1.4}
\begin{tabular}{|p{3.1cm}|p{5.4cm}|p{2.6cm}|}
\hline
\textbf{Category} & \textbf{Condition} & \textbf{Applies to} \\\hline
Offensive jihad & Requires an imam's presence; not permitted to a lone authority acting alone & Neither brother's case \\\hline
Defensive jihad & Obligatory on all when Islamic land/people/property attacked & Closer to al-Husayn \\\hline
Jihad-equivalent (self-defense) & Fighting to save one's own life; scales with what's at stake & al-Husayn at Karbala \\\hline
Fighting Muslim rebels (Qur'an 49:9) & Reconcile first; fight only the side that refuses & Available, not obligatory, to al-Hassan \\\hline
Legitimate truce & Only if too weak to prevail, need time, or truce may draw the other side toward Islam; fixed term only & al-Hassan's treaty \\\hline
\end{tabular}
\vspace{0.3cm}
\delivernote{Walk the room through this table before saying a word about the actual treaty -- that ordering is deliberate, so nobody judges the treaty by gut feeling before knowing the legal categories.}
\end{frame}

\begin{frame}{Hudaybiyyah: The Direct Precedent}
\begin{itemize}[itemsep=5pt]
  \item Apparently one-sided 10-year truce: Muslim defectors returned, no reciprocal clause
  \item 'Umar objected; the Prophet accepted anyway
  \item Within under two years: more conversions than the previous two decades of open conflict
\end{itemize}
\bigquote{Leave me and seek someone else... I would rather stay a minister than to become a chief.}{'Ali -- a parallel restraint, decades earlier}
\end{frame}

% ---- Jihad & Truce decision tree -------------------------------------------
% Four full-bleed 16:9 figures generated from figures/jihad_truce_tree.json by
% scripts/build_jihad_truce_tree.py (regenerate; do not edit the PDFs by hand).
% They replace the earlier two text-column frames "Follow the Yes/No Branches" and
% "Two Historical Paths Overlaid on the Framework"; their caveats are kept in the notes.
\newcommand{\treeframe}[2]{%
  {\setbeamertemplate{background}{\includegraphics[width=\paperwidth,height=\paperheight,keepaspectratio]{figures/jihad-truce-tree-#1.pdf}}%
   \begin{frame}[plain]
   \delivernote{#2}
   \end{frame}}%
}

\treeframe{clean}{Walk the tree left to right before any Imam appears. Stage 1 classifies the conflict: defence, fighting rebels, or enjoining good. Stage 2 asks whether a truce is lawful and wise. This is Mutahhari's source-based teaching framework, not a stand-alone ruling for modern conflicts, and no single branch settles the historical comparison. Ask the room to predict which exits each Imam will take.}

\treeframe{overlay}{Both Imams pass through the same boxes. Green is al-Hassan and red is al-Husayn; they enter the truce box together and leave in opposite directions. Read the scorecard: on the three conditions for rising, al-Hassan scores 0 of 3 and al-Husayn 3 of 3. Same principles do not mean identical actions. Mutahhari's argument is cumulative: role, capacity, demanded allegiance, public conditions, treaty terms and foreseeable consequences must be read together. The colours mark each route; poison and bloodshed are historical aftermath, not causal outputs of the tree.}

\treeframe{hassan}{Follow the green line. Applicable: rebellion against the sitting caliph; weak and divided support; peace as an available option; defined treaty terms. Not implied: a pledge of allegiance to every act of Mu'awiyah. Aftermath: the tradition reports that al-Hasan was later poisoned. This is not a causal output of the tree or proof that the treaty caused his death.}

\treeframe{husayn}{Follow the red line. Applicable: direct demand for allegiance to Yazid; the ruler's record; public invitations; responsibility to oppose wrongdoing. Not a simple match: his case is not settled by classifying it as offensive jihad or applying truce rules alone. Aftermath: the stand culminated in Karbala's bloodshed and martyrdom; this is historical aftermath, not a universal algorithm.}

\begin{frame}{Same Principles, Different Historical Conditions}
\small
\renewcommand{\arraystretch}{1.25}
\begin{tabular}{|p{2.3cm}|p{4.4cm}|p{4.4cm}|}\hline
\textbf{Question} & \textbf{Imam al-Hasan: treaty} & \textbf{Imam al-Husayn: qiyam}\\\hline
Political role & Sitting caliph facing rebellion; his killing would also mean the killing of the caliph in office. & Not the sitting caliph; refusal became public protest against Yazid's rule.\\\hline
Military and social conditions & Mutahhari describes divided forces, unreliable support, and the likely cost of a prolonged war. & The earlier treaty situation had changed; the later case included a demand for allegiance and public invitations from Kufah.\\\hline
Moral/political trigger & The same immediate demand for an oath of allegiance was not made of al-Hasan. & Yazid's demand for allegiance, public opposition to wrongdoing, and the claim created by Kufah's invitations.\\\hline
Terms / purpose & Treaty clauses aimed to constrain governance, succession, public conduct, and reprisals against supporters. & The stand was not simply a troop-count calculation; Mutahhari stresses the ruler's record and public responsibility.\\\hline
\end{tabular}
\vspace{0.2cm}
\begin{block}{Read the argument cumulatively}
No single row is a universal decision rule. Mutahhari's argument depends on the combination of role, military capacity, allegiance, public conditions, and foreseeable consequences.
\end{block}
\end{frame}

\begin{frame}{Five Contrasts: al-Hassan vs. al-Husayn}
\centering\scriptsize
\renewcommand{\arraystretch}{1.4}
\begin{tabular}{|l|p{4.7cm}|p{4.7cm}|}\hline
&\textbf{al-Hassan}&\textbf{al-Husayn}\\\hline
Position & Sitting caliph -- death parallels 'Uthman's assassination & No office -- death is honor, not institutional collapse\\\hline
Military balance & Iraq (30k--100k max) vs. Mu'awiyah's 150k; Kufah exhausted, bribed & Different balance a generation later\\\hline
Allegiance demand & None -- treaty explicitly excludes any oath & Yazid's explicit personal demand -- the trigger\\\hline
Kufah's mood & Exhausted, hostile -- wore armor even to lead prayer & 18{,}000 letters -- a mass mandate history would fault him for ignoring\\\hline
Ruler's record & Mu'awiyah's tyranny not yet proven; offered signed blank paper & Two decades of documented injustice by Yazid's time\\\hline
\end{tabular}
\delivernote{No single row is decisive alone -- the strength of the argument is that on every axis Mutahhari can find, the two situations genuinely diverge.}
\end{frame}

\begin{frame}{What the Treaty Actually Secured}
\begin{itemize}[itemsep=4pt]
  \item Governance to Mu'awiyah \emph{conditioned} on ruling by Qur'an and Sunnah
  \item Succession reverts to al-Hassan, then al-Husayn -- explicitly temporary
  \item End to public cursing of 'Ali in Syria
  \item Treasury funds for al-Hassan's household and Siffin/Jamal war-widows
  \item General amnesty for 'Ali's former supporters
\end{itemize}
\bigquote{}{Mu'awiyah later broke essentially every one of these terms -- which is what exposed his true character, not the peace itself.}
\end{frame}

\qaframe{al-Hassan's Pacifism}{%
  \qa{Lay out all five contrasts. Does any single one, alone, justify al-Hassan's peace -- or does the argument only work cumulatively?}
     {No single factor is decisive alone. Removing any one wouldn't flip the verdict, since the others still stand independently. The strength is that on every axis Mutahhari can find, the two brothers' situations genuinely diverge.}
  \qa{If Hudaybiyyah is the model, what is the risk of treating every costly-looking peace as automatically wise in hindsight?}
     {Survivorship bias -- Hudaybiyyah is remembered as vindicated because it worked. The specific jurisprudential conditions (weakness, need for time, likelihood of drawing the other side toward Islam) have to actually hold in a given case, not just be assumed by loose analogy to a famous precedent.}
  \qa{Play devil's advocate: what is the strongest counter-argument a critic of al-Hassan could still make?}
     {The five contrasts explain why war would have been strategically unwise, but don't fully settle whether al-Hassan had a moral duty to resist regardless of the odds. Mutahhari's implicit answer: obligation scales with genuine possibility and manifest cause, not a blanket principle that surrender is always shameful.}
}
